\section{Verification and open questions}\label{sec:verify}

We computed in Python, in double precision unless stated otherwise. For
the elephant walk, a forward dynamic program for the law of $M_n$ and a
backward recursion for $V_s$, folded by the symmetry $V_s(b)=V_s(s-b)$,
agree to $1.6\cdot10^{-14}$ at $n=12800$. They also agree with exact
rational arithmetic for $n\le60$ and with 40-digit arithmetic at
$n=400$. A recursion for Taylor coefficients in $\varepsilon$ gives $c_1$
and $k_2(n)$. For the Markov crossing we count words by their number of
switches, and this reproduces the values of $p_n$ in~\cite{paperA}.
Section~\ref{sec:shape} uses exact recursions for $f_n$ and $\CP_n$.
Truncating the chain at a level $m_0$ does not change $f_n(m)$ for
$m\le m_0$, since $M_t$ never decreases. For heavy tails we compared 3
methods. These are the two-renewal formula of
Proposition~\ref{prop:tworenewal}, a recursion on the first run with exact
Hurwitz zeta tails, and a sum over all $2^{n-1}$ run compositions in
30-digit arithmetic. All 3 agree to $7\cdot10^{-16}$ for $n\le14$, and the
first 2 agree to $10^{-14}$ up to $n=3200$. For aging,
a forward recursion with nonnegative terms and a transfer matrix
evaluated at roots of unity agree to at least 12 digits on the crossings,
and small cases were done in exact rationals. The constants of the Landau
density come from quadrature with 25 digits. Programs written separately
reproduced the main numbers, including values of $x_n$, of $D(n)$, of
$R_n$ and of $c_n$.

There is one check script for each part of the paper.
\texttt{verify\_elephant\_crossing.py} covers Section~\ref{sec:erw},
Table~\ref{tab:erw-crossing} and Appendix~\ref{app:erw}.
\texttt{verify\_elephant\_shape.py} covers Section~\ref{sec:shape} and
Appendix~\ref{app:shape}. \texttt{verify\_heavy\_runs.py} covers
Section~\ref{sec:heavy}, Table~\ref{tab:heavy} and
Appendix~\ref{app:heavy}, and \texttt{verify\_aging.py} covers
Remark~\ref{rem:aging} and Appendix~\ref{app:aging}. Each script
recomputes the numbers printed in its part and compares them with the
printed values. Where 2 independent methods were used when a number was
first found, the script runs both. \texttt{verify\_all.py} runs the 4
scripts. The quick run takes a few minutes, and the full run adds the
largest cells and takes less than an hour. All 297 checks of the full run
pass. The scripts need numpy, scipy and mpmath, and their output is in
\texttt{paperD/data}.

Many of the exact and finite statements are also proved in Lean~4, in
\texttt{lean/PaperD} in the Problem 131 folder. For the elephant walk
these are Lemma~\ref{lem:erw-reduction}, Proposition~\ref{prop:mlr}, the
value $C_n'(0)=4/(n+2)$ of Proposition~\ref{prop:first-mutation}, the 2
sums of Lemma~\ref{lem:subtree} (for the urn that gives the subtree
sizes), Lemma~\ref{lem:unimodal} for $k_0=1$, Theorem~\ref{thm:dip},
Corollary~\ref{cor:dip} and Lemma~\ref{lem:B-D}. For heavy tails they are
Proposition~\ref{prop:edge}, the identities for $e(\alpha)$ and
$K_\alpha$ behind Corollary~\ref{cor:golden}, and the value
$\ell_\alpha(\frac12)$ in Theorem~\ref{thm:lamperti}. For aging they are
Propositions~\ref{prop:D-paths} and~\ref{prop:D-uniform},
Corollary~\ref{cor:D-fixed}, Lemma~\ref{lem:D-average} and the finite
part of Theorem~\ref{thm:D-bracket}. The file \texttt{MAP.md} in that
folder lists the Lean theorems. The asymptotic statements are not in
Lean.

The proofs do not use these computations, with 2 exceptions. The
existence of $\kappa_L$ in Theorem~\ref{thm:landau}(c) uses the sign of
$f_L''(\lambda_0)$, which we computed but did not certify. Also some
constants in the proofs are values of explicit series and integrals, which
we evaluated numerically. Conjecture~\ref{conj:convex}, the checked ranges
after Corollary~\ref{cor:dip}, and the numerical parts of
Remarks~\ref{rem:refined}, \ref{rem:fresh}, \ref{rem:crossover}
and~\ref{rem:heat} rest on computation alone.

Before computing a new case we wrote down, in the ledgers
\texttt{ledger\_D1.md} to \texttt{ledger\_D4.md} in \texttt{paperD/data},
a prediction and a condition that would refute it. Outcomes are recorded
below each prediction, refuted predictions were kept, and observations
made after a computation are marked as such. Most predictions survived,
and 8 failed. The law $f_n$ given the first step is not convex between
$2m^*$ and $n-2m^*$ (item D2-P9), which led us to
Conjecture~\ref{conj:convex} for the full law. At $n=10^6$ and $x=320$
the mode of $f_n$ is 3 below the mode of $\CP_n$, where we allowed 2
(D2-P11). The endpoint dip already holds at $n=4$ and $5$, where we
predicted that it fails (D2-P12). Two scaling laws for the finite-$n$
shift between these 2 modes were both wrong (D2-P15). The form of
Conjecture~\ref{conj:convex} on the closed range $[2m^*,n-2m^*]$ fails
for 18 values of $n\le51$, which a separate check of $10\le n\le99$, also
written down in advance, found. For heavy tails, one tolerance was too
small because of an arithmetic slip in the prediction (D3-P2), and one
predicted value of $n^*(1.60)$ was off by 2 because we misread a rounded
number (D3-V1). For aging, we predicted $nC_n=d_c(1+(c-1)/n+O(n^{-2}))$,
with $d_c$ the density of $\mathrm{Beta}(c,c)$ at $\frac12$, and at
$c=0.3$ the remainder is larger (D4-P7). None of the failures concerns a
statement that we call a theorem.

Several questions remain open. For the elephant walk, the refined
crossing of Remark~\ref{rem:refined} needs an upper bound
$\log C_n\le\log(\varepsilon c_1)+\varepsilon k_2(n)
+O(\varepsilon^2(\log n)^k)$, which we do not have. We have no explicit
$n$ beyond which $g_n>x_n$ (Theorem~\ref{thm:erw-crossing}). Evaluating
the lower bound of Theorem~\ref{thm:erw-bounds}(c) at $g_n$ suggests that
$n\ge2^{17}$ is enough, and $x_n<g_n$ at every $n$ in
Table~\ref{tab:erw-crossing}, but we have not proved it.
Theorem~\ref{thm:shape} says nothing in a window around the center, and
closing it needs second differences of the law, as in
Conjecture~\ref{conj:convex}. Also open are the tail of
Theorem~\ref{thm:profile} for $m$ between $x\log x$ and $x\log^2x$, a
sharp upper bound on the mirror term for $m$ of order $n$, and effective
constants in Theorems~\ref{thm:landau}(c) and~\ref{thm:shape}, starting
with a certified value of $\kappa_L$. For heavy tails we have no rate of
convergence in Theorem~\ref{thm:center}. The monotonicity of $R_n$, and
the fact that $R_n-1$ changes sign only once, at $n^*(\alpha)$, are only
observed on the computed range (Remark~\ref{rem:crossover}). The case
$\alpha=1$ is only heuristic (Remark~\ref{rem:fresh}). We have not
studied the whole law for $1<\alpha<2$, for example whether the end atoms
are isolated maxima. For $\alpha<1$ we know the shape of $\ell_\alpha$
only near $y=\frac12$, and we do not know whether the center bin at
finite $n$ is a local minimum for $\alpha<\alpha_c$ and a local maximum
for $\alpha>\alpha_c$. For aging, we checked that at $c=c_n$ the center
is the unique interior minimum and bins $1$ and $n-1$ are the largest,
for every even $n$ from 4 to 1000 and for $n=1600,3200,6400$, but we
have not proved this.

\paragraph{Acknowledgment.}
The Lean code for this paper was generated by AI. AI was also used to
revise some of the writing and to help cross-verify the code.

\paragraph{Data and code availability.}
The manuscript source, verification scripts, recorded results and the
prediction ledgers are in the Problem 131 folder of
\url{https://github.com/apemm/Kagey-Problems}, under \texttt{paperD}. The
Lean files are in \texttt{lean/PaperD} in the same folder.
